\section{The smooth cutoff obstruction in one-sided printed norms}
\label{sec:printed-equivariant-cutoff-obstruction}

\paragraph{Status.}
The complete post-writing author reading and two independent
whole-proof readings have passed.
It concerns an explicitly specified smooth target.
No Hilbert-only assertion at the untwisted line is inferred.

Let \(K\) be a fixed number field and let \(\eta\) be a fixed
primitive finite-order Hecke character, including the principal
character.  Write \(\Xi_\eta\) for its entire pole-removed
completion, and \(m_\rho\) for the multiplicity of its
nontrivial zero \(\rho\).
Let \(\mathcal E\) have the full exponential-Schwartz
topology used in MT-M18 and MA-PI-005, and put
\[
 \begin{gathered}
  \mathcal I_\eta^{\,c}
   =\{h\in\mathcal E:
       (\mathcal Mh)^{(j)}(\rho-\tfrac12)=0
       \text{ for all }\rho,\ 0\le j<m_\rho\},\\
  \mathcal Mh(z)=\int_{\mathbb R}e^{zx}h(x)\,dx,\qquad
  (\tau_yh)(x)=h(x-y),\\
  H^{\rm one}_{b,s}
   =L^2\bigl(\mathbb R,e^{2bx}(1+x^2)^{s/2}\,dx\bigr),
       \quad b\in\mathbb R,\ s\ge0 .
 \end{gathered}
 \tag{MA-PI.170}
\]
The first line is the actual centered smooth ideal by
the full-jet and theta-image identifications in MT-M18/MT-M24.
For the argument it would suffice to assume the displayed
vanishing conditions; no equality of differently completed
ideals is needed.

\paragraph{MA-PI-007 (one-sided norms and fixed Hecke smooth sectors).}
Let \(T:\mathcal D(\mathbb R)\to\mathcal I_\eta^{\,c}\)
be continuous in the full \(\mathcal E\) topology and
translation-equivariant.  Fix \(b\in\mathbb R\) and \(s\ge0\).
Suppose there is a fixed \(R\ge0\) such that, whenever
\(\sup\operatorname{supp}f+R<\inf\operatorname{supp}g\),
\[
 \big\langle T\tau_{-L}f,T\tau_{-L}g
                      \big\rangle_{H^{\rm one}_{b,s}}=0
       \quad\text{for every sufficiently large }L .
 \tag{MA-PI.171}
\]
The threshold may depend on \(f,g\).
Then \(T=0\).
The same assertion holds with domain
\(\mathcal D_0=\{f\in\mathcal D:\int e^{-x/2}f(x)\,dx=0\}\)
and its inherited topology.

The assertion also holds for the symmetric norm
\(H_{a,s}=L^2(e^{2a|x|}(1+x^2)^{s/2}\,dx)\), \(a>0\),
with (MA-PI.171) interpreted in that norm.
Thus the Hecke-sector extension applies to both exact
metrics, while the printed one-sided case permits every
real \(b\), including \(b=0\).
These are smooth-target statements, without a narrower
zero-location hypothesis, A, M, or a simplicity hypothesis.

\begin{proof}
\emph{1. The exact limiting Gram form.}
For \(h,k\in H^{\rm one}_{b,s}\), changing variables gives
\[
 \begin{split}
 &e^{2bL}L^{-s}
       \langle\tau_{-L}h,\tau_{-L}k\rangle_{H^{\rm one}_{b,s}}\\
 &\quad=\int_{\mathbb R}h(x)\overline{k(x)}e^{2bx}
       \frac{(1+(x-L)^2)^{s/2}}{L^s}\,dx\\
 &\quad\longrightarrow
       \int_{\mathbb R}h(x)\overline{k(x)}e^{2bx}\,dx .
 \end{split}
 \tag{MA-PI.172}
\]
For \(L\ge1\), the polynomial ratio is bounded by
\(C_s(1+x^2)^{s/2}\) and tends pointwise to one.
Cauchy--Schwarz in the original displayed norm proves
the required domination.  Consequently (MA-PI.171)
gives separated-support orthogonality for the weight
\(e^{2bx}\).

For the symmetric norm use the already proved exact
limit (MA-PI.154): its normalization is
\(e^{-2aL}L^{-s}\), and its limiting weight is \(e^{-2ax}\).
Thus that case reduces to the following argument with
\(b=-a\).  No orthogonal projection has been transported
by a merely equivalent norm.

\emph{2. A compactly supported Gram distribution.}
Define
\[
 \begin{gathered}
   Af(x)=e^{bx}T(e^{-b\cdot}f)(x),\qquad
   \widehat f(t)=\int_{\mathbb R}e^{itx}f(x)\,dx,\\
   \widehat{Af}(t)=M_b(t)\widehat f(t),\qquad
   \widehat k(t)=|M_b(t)|^2,\qquad
   \operatorname{supp}k\subset[-R,R].
 \end{gathered}
 \tag{MA-PI.173}
\]
Here the last three objects exist by the same actual
construction as in the proof of MA-PI-005.
For clarity, the construction uses only the following
facts, all valid for this arbitrary real \(b\).
The inclusion \(\mathcal E\to L^2(e^{2bx}dx)\) is continuous,
so \(A:\mathcal D\to L^2\) is continuous.
The exponential factors cancel under translation, making
\(A\) translation-equivariant.  Its separated-support
orthogonality follows from Step~1 because multiplication
by \(e^{-bx}\) preserves supports.

A unit-interval partition of unity and the continuity
estimate on one fixed compact interval extend \(A\)
to a continuous map \(\mathcal S\to L^2\): translating
each partition piece back to that interval gives the
same finite derivative bound, and two extra Schwartz
powers make the sum of \(L^2\) norms convergent.
The convolution identities
\(A(f*g)=(Af)*g=f*(Ag)\) then hold in \(L^2\).
Choosing a Gaussian with nowhere-zero Fourier transform
defines \(M_b\) as the quotient of its image transform
by its own transform.  It is locally \(L^2\).
Modulating the Gaussian and applying the Schwartz
continuity estimate bounds the \(L^2\) mass of \(M_b\)
on each unit interval by a polynomial in its center.
Thus \(|M_b|^2\) is a tempered distribution.
Plancherel identifies its inverse Fourier distribution
\(k\) with the Gram form.  Testing that form on two
separated compact supports, and taking a compact
approximate identity in one factor, makes \(k\)
vanish outside \([-R,R]\).  Interchanging the two
inputs supplies both components of the complement.

The Fourier transform \(\mathfrak P\) of this compact
distribution therefore has an entire extension with
\[
 \begin{aligned}
 |\mathfrak P(w)|
   &\le C(1+|w|)^N e^{R_1|\operatorname{Im}w|},\\
 m(z)\overline{m(2b-\overline z)}
   &=\mathfrak P\bigl(-i(z-b)\bigr).
 \end{aligned}
 \tag{MA-PI.174}
\]
Here \(R_1\) may be any fixed number larger than \(R\).
The bound follows directly by applying the finite-order
compact-distribution estimate to \(x\mapsto e^{iwx}\).
If \(\mathfrak P\ne0\), Jensen's formula about a point
where it is nonzero gives \(O(T)\) zeros in \(|w|\le T\),
counted with multiplicity.

\emph{3. The entire multiplier and the full Hecke divisor.}
To justify the second identity in (MA-PI.174), apply the
Mellin transform to
\(T(f*g)=(Tf)*g=f*(Tg)\) for compact tests.
The transforms of the outputs are entire because their
values lie in \(\mathcal E\).
At every \(z_0\in\mathbb C\) a compact test has nonzero
Mellin value there.  The quotients
\(\mathcal M(Tf)/\mathcal Mf\) patch on these open
neighborhoods to an entire function \(m\), with
\(\mathcal M(Tf)=m\,\mathcal Mf\) for every compact test.
On \(z=b+it\), conjugation in the definition of \(A\)
gives \(M_b(t)=m(b+it)\) almost everywhere.
The equality in (MA-PI.174) holds first on this line,
then everywhere by the identity theorem.
The sign of its argument is determined by the
plus-Fourier convention in (MA-PI.173).

At a zero \(\rho\), choose an input whose Mellin
transform is nonzero at \(\rho-\tfrac12\).
The full-jet condition in (MA-PI.170) implies
\(\operatorname{ord}_{\rho-1/2}m\ge m_\rho\).
Every such multiplicity therefore occurs among the
zeros of \(\mathfrak P\) at
\(-i(\rho-\tfrac12-b)\).
For this fixed primitive finite-order character the
classical unconditional count is
\[
 \sum_{\substack{\Xi_\eta(\rho)=0\\|\operatorname{Im}\rho|\le T}}
       m_\rho
      =\frac{[K:\mathbb Q]}{\pi}T\log T+O_{K,\eta}(T).
 \tag{MA-PI.175}
\]
The precise primary input is Greni\'e--Molteni (2019),
Appendix~A, Theorem~A.1, author p.~26.
Its definition counts both ordinate signs and all
multiplicities, and its statement explicitly includes
every primitive finite-order Hecke character.
The nontrivial zeros lie in the classical critical strip,
so their transformed points have modulus at most
\(T+|b|+1\).  This contradicts the \(O(T)\) zero bound
unless \(\mathfrak P=0\).  Hence \(M_b=0\) almost
everywhere, \(A=0\), and \(T=0\).
No placement inside a narrower strip has been used.

\emph{4. The moment domain.}
For a map defined on \(\mathcal D_0\), precompose with
\(D-\tfrac12:\mathcal D\to\mathcal D_0\).
This differential map is continuous, support-preserving
and translation-equivariant, and it is onto by the
compact primitive displayed in MA-PI-005.
The preceding proof applies to this composite.
Surjectivity gives the asserted conclusion on
\(\mathcal D_0\).
\end{proof}

\paragraph{The exact rational printed realization.}
Let \(0<A<1\) now denote Masullo's analytic twist parameter,
and put \(b=A-\tfrac12\).
In the compact-trivial sector of the original
\(W^{s,2}(C_{\mathbb Q})\), let
\(m_1=\mu_1(C_{\mathbb Q}^1)>0\).
The exact unitary of MA-PI-002 specializes to
\[
 (U_bh)(\nu(e^x)k)
      =m_1^{-1/2}e^{bx}h(x),\qquad
 U_b:H^{\rm one}_{b,s}
       \overset{\cong}{\longrightarrow}
             W^{s,2}(C_{\mathbb Q})(\mathbf1).
 \tag{MA-PI.176}
\]
The compact Haar mass and the actual product measure
are retained.  The source twist is \(|g|^A\), whereas
the theta normalization supplies the factor
\(|g|^{1/2}\); their quotient is exactly \(e^{bx}\).
This is an isometry, including \(A=1/2\), \(b=0\).

Let \(\mathfrak T:\mathcal D_0\to
W^{s,2}(C_{\mathbb Q})(\mathbf1)\) be a candidate
transport on the genuine global PW chart (MA-PI.162).
Assume explicitly that \(T=U_b^{-1}\mathfrak T\) is
continuous into the full centered smooth theta ideal
and translation-equivariant.
These are the centered smooth-source hypotheses;
mere continuity into the Hilbert space is not substituted.
Let \(Q^{\rm src}_{Y,A,s}\) be the original printed
orthogonal PW projection.
If the exact physical-cutoff transport supplies
\[
 \begin{array}{ll}
 Q^{\rm src}_{Y,A,s}\mathfrak T v=\mathfrak T v,
       &\operatorname{supp}v\subset[c+\log2,\infty),\\
 Q^{\rm src}_{Y,A,s}\mathfrak T v=0,
       &\operatorname{supp}v\subset(-\infty,c-\log2],
 \end{array}
 \qquad c=-\log Y
 \tag{MA-PI.177}
\]
for all sufficiently large \(Y\), then \(\mathfrak T=0\).
Indeed the exact unitary transfers these orthogonal
plateau identities to \(H^{\rm one}_{b,s}\).
For two source supports separated by more than
\(\log4\), choose \(c\) between the plateaus.
After both supports are translated far left, \(Y\)
is sufficiently large.  One projector fixes one image
and kills the other, so (MA-PI.171) holds with
\(R=\log4\).  MA-PI-007 applies.

The physical chart, its two moment conditions and
its dilation law are exactly those checked in MA-PI-005,
using the original global PW source on Masullo
pp.~354--355, Definition~7.3.2.3 and (7.3.2.12)--(7.3.2.17).
The cutoff is the Fourier ball in p.~516, Step~II,
with the eventual plateaus justified by the p.~519
thin-transition assumptions.
The metric and twist use pp.~399--400, (9.1.2.3),
pp.~434--435, (10.1.1.1)--(10.1.1.3), and
pp.~488--489, (11.2.1.7)--(11.2.1.8), as in MA-PI-002.
In particular the conclusion holds at the actual
\(A=1/2\) specified on p.~515.
The theorem concerns the stated chart and transport
hypotheses.  It does not identify that chart with
the coordinate-flat local source.

\paragraph{Hecke-sector and framework scope.}
For a nontrivial Hecke character, MA-PI-007 is already
an abstract smooth-sector theorem with its actual full
completed divisor.  An application to physical source
data must additionally verify its own chart, finite
Fourier factors, dilation covariance and two plateaus.
Those are not supplied by a formal change of character.
Nor does the one-sided proof assert the Hilbert-only
variant of MA-PI-006 at \(b=0\).

The inherited SGA~1, V~\S7, and Deligne
\S\S6.1--6.3 locators place a finite-order character
in its finite-monodromy rank-one realization.
Its local Euler determinants are unchanged.
The parameters \(A,b\) above are analytic norm
characters; half-centering is not identified with an
integral Tate twist.
The finite-length zero-jet modules enter exactly
through their multiplicities.
In the inherited cote~1 BR-TR-01/02 framework,
scans079--080 and111--112, this obstruction occurs
before an intertwiner or nuclear trace pairing can
be assigned.  It makes no positivity or motivic
weight inference and does not settle FT002.

\paragraph{Primary arithmetic reference and reading scope.}
Greni\'e, L., \& Molteni, G. (2019).
An explicit Chebotarev density theorem under GRH.
\emph{Journal of Number Theory, 200}, 441--485.
\url{https://doi.org/10.1016/j.jnt.2018.12.005}.
Only its unconditional Appendix~A zero-count theorem
is invoked.  The primitive finite-order setup,
both-sign multiplicity convention and Theorem~A.1
statement on author pp.~25--26, and the completion
definitions (2.5)--(2.8) on p.~5, were freshly read.
Only the classical leading-order count is used;
none of the optimized numerical error constants
in Theorem~A.1 is invoked.
No fresh complete appendix proof audit is claimed.
